%%%%%%%%%%%%%%%%%%%%%%%%%%%%%%%%%%%%%%%%%%%%%%%%%%%%%%%%%%%%%%%%%%%%%
%% sec_intro.tex
%% Section 1: Introduction
%% RMG Periodic RT-TDDFT -- Paper 2
%%%%%%%%%%%%%%%%%%%%%%%%%%%%%%%%%%%%%%%%%%%%%%%%%%%%%%%%%%%%%%%%%%%%%

\section{Introduction}
\label{sec:intro}

Real-time time-dependent density functional theory (RT-TDDFT) provides
a direct first-principles route for following electronic response after
an applied perturbation.  By propagating the electronic degrees of
freedom in time, RT-TDDFT gives access to optical response,
charge-transfer dynamics, beam-induced electronic excitations,
plasmonic excitation and energy transfer, nonlinear response, and
ultrafast electron dynamics within a single initial-value
framework.\cite{tddft-fundamentals,chem-rev-Li,bertsch-96,lvn-ct-2012,JCTC-2020,tao-2021}
For extended materials, low-dimensional systems, magnetic compounds, and
defects in solids, the same real-time perspective is attractive because
the relevant response may depend simultaneously on crystal momentum,
spin polarization, local defect states, and long-range periodic
environment.  These requirements motivate scalable RT-TDDFT
implementations that can treat periodic boundary conditions and large
supercells without abandoning a first-principles electronic-structure
description.

Periodic systems require special care in the representation of a
homogeneous optical perturbation.  In finite systems, the electric-dipole
interaction is often written in the length gauge through the position
operator.  Under periodic boundary conditions, however, the ordinary
position operator is not a lattice-periodic observable, and the
macroscopic polarization of an extended insulator is defined through
Berry-phase or related periodic formulations rather than through a
simple dipole
moment.\cite{Resta-polarization-1994,KingSmith-Vanderbilt-polarization-1993,periodic-operator}
A natural periodic formulation of the long-wavelength optical response
therefore uses a spatially homogeneous vector potential and evaluates
the induced current response.\cite{Luber-rt-tddft,Hanasaki-2023}
The homogeneous vector potential represents the optical
$\mathbf q\rightarrow0$ limit; the restriction to vertical transitions
comes from the long-wavelength approximation, not from the velocity
gauge itself.

Table~\ref{tab:external-field-map} places this homogeneous optical
limit in the broader landscape of external-field representations used
in time-dependent electronic-response theory.  The table is intended as
a conceptual map for the present discussion, not as a list of all
capabilities implemented in \texttt{RMG}.  In Coulomb gauge,
$\nabla\cdot\mathbf A=0$, the vector potential is transverse; after the
longitudinal/transverse decomposition of the electric field,
$-\nabla\Phi$ gives the longitudinal electric field, while
$-(1/c)\partial\mathbf A/\partial t$ gives the transverse electric field
and $\nabla\times\mathbf A$ gives the magnetic field.  The present work
addresses the homogeneous periodic optical-response row of
Table~\ref{tab:external-field-map}.  The finite-$\mathbf q$, magnetic,
and multipolar rows are included only as broader perspective for related
TDDFT response problems.\cite{Rubio-2002,Luber-rt-tddft}

\begin{table}[t]
\scriptsize
\setlength{\tabcolsep}{3pt}
\hbadness=10000
\caption{Summary of representative external-field perturbations relevant
to time-dependent electronic-response calculations.  The table is
intended as a conceptual map.  Among the perturbations summarized here,
the present work treats the homogeneous periodic optical response.}
\label{tab:external-field-map}
\centering
\begin{tabular}{@{}p{0.16\linewidth}p{0.24\linewidth}p{0.17\linewidth}p{0.20\linewidth}p{0.12\linewidth}@{}}
\hline
Physical probe / response
& Representative perturbation
& Field character
& Typical applications
& Relation to this work \\
\hline
Homogeneous optical electric field
& Length-gauge dipole coupling for finite systems; homogeneous
  $\mathbf A(t)$ velocity-gauge representation for periodic systems
& Long-wavelength $\mathbf q\rightarrow0$ optical limit
& Optical response, ultrafast laser-driven electron dynamics,
  pump/probe-type response
& Present work \\[0.6em]
Finite-$\mathbf q$ longitudinal scalar perturbation
& $\Phi_{\mathbf q}(t)e^{\imath\mathbf q\cdot\mathbf r}$
& Longitudinal electric field / density response
& Finite-$\mathbf q$ density response, EELS/loss spectra,
  plasmon dispersion
& Broader perspective; not implemented here \\[0.6em]
Finite-$\mathbf q$ transverse electromagnetic perturbation
& $\mathbf A_{\mathbf q}(t)e^{\imath\mathbf q\cdot\mathbf r}$,
  with $\mathbf q\cdot\mathbf A_{\mathbf q}=0$ in Coulomb gauge
& Transverse electromagnetic response with associated magnetic field
& Finite-$\mathbf q$ transverse optical or radiative response beyond
  the homogeneous optical limit
& Broader perspective; not implemented here \\[0.6em]
Magnetic / spin perturbation
& Time-dependent magnetic field, Zeeman coupling, or spin-dependent
  perturbation
& Magnetic or spin response
& Spin dynamics, magnetic resonance, magnetic excitations, magnons in
  appropriate noncollinear spin-dynamical frameworks
& Not presently implemented as such \\[0.6em]
Multipolar / chiroptical perturbations
& Electric quadrupole, magnetic dipole, field gradients, or mixed
  electric/magnetic response
& Beyond electric-dipole response
& Beyond-dipole spectroscopy, ECD, MCD
& Broader perspective; not implemented here \\
\hline
\end{tabular}
\end{table}

RT-TDDFT for extended systems has been developed in a variety of
numerical representations.  Real-space implementations such as
\texttt{Octopus} provide time propagation for finite and extended
systems and serve here as the independent reference for cross-code
validation.\cite{octopus-2020}  Plane-wave and related periodic
approaches, including \texttt{Qbox}/\texttt{Qball} developments and
\texttt{CP2K}, have demonstrated real-time response calculations under
periodic boundary
conditions.\cite{qball-2017,qball-II-2021,Luber-rt-tddft,cp2k-2020}
All-electron and atom-centered formulations provide complementary
routes for crystalline and molecular systems, including recent periodic
RT-TDDFT implementations in \texttt{exciting}, numeric atom-centered
basis frameworks, and PySCF.\cite{exciting-code-2021,Hekele-2021,Hanasaki-2023}
Finite-element approaches provide another real-space route to large
systems.\cite{gavini-rt-tddft}  In a related approximate direction,
Wong and co-workers have developed a velocity-gauge real-time TD-DFTB
implementation for periodic condensed-matter systems based on the DFTB+
framework.\cite{Xu-Wong-RT-TDDFTB-2023,DFTB+}
These developments show that periodic real-time response is an active
area with multiple complementary implementations; the present work adds
the particular combination of periodic current response, collinear spin
resolution, finite active-space density-matrix propagation, and
real-space multigrid scalability in \texttt{RMG}.

Building on our previous finite-system \texttt{RMG} RT-TDDFT
implementation, which demonstrated real-time simulations of systems
containing up to approximately 24,000 electrons on massively parallel
architectures,\cite{jakowski-rmg-tddft-2025} the present work extends
the framework to periodic boundary conditions and collinear
spin-polarized response.  This extension is not only a change of
boundary conditions.  Periodic materials require Bloch-orbital and
k-point resolution, a current-based optical observable, and careful
treatment of nonlocal pseudopotential contributions to the current
operator.  Spin-polarized systems further require separate density
matrices and self-consistent densities for the two collinear spin
channels.  The implementation also builds on the real-space multigrid
infrastructure of \texttt{RMG}, which is designed for large-scale
parallel calculations on modern CPU and GPU architectures.\cite{rmg-dft-2024}

The main contributions of this paper are as follows.  We formulate
periodic RT-TDDFT for homogeneous optical response in the velocity
gauge, with the macroscopic response expressed through the current.  We
derive the periodic current operator, including the contribution from
norm-conserving nonlocal pseudopotentials.  We extend the density-matrix
propagation scheme to finite active spaces resolved by Bloch wave vector
and collinear spin channel.  We then describe the corresponding
implementation in \texttt{RMG}, including k-point, spin-channel,
real-space, and distributed-matrix parallel organization.

We validate the periodic and spin-polarized implementation by comparing
current-derived optical spectra from \texttt{RMG} with independent
RT-TDDFT calculations for a sequence of increasingly demanding systems:
a one-dimensional periodic hydrogen-dimer reference, pristine monolayer
hBN, bulk Si, and spin-polarized \ch{CrI3}.  These comparisons test the
principal spectral structure rather than absolute optical intensity.  As
a large-scale scientific demonstration, we apply the method to the
negatively charged nitrogen-vacancy center in diamond, a spin-polarized
defect whose optical response probes local defect orbitals,
polarization selection rules, spin-channel selectivity, active-space
convergence, and defect-concentration effects.\cite{NV-Doherty-2013,NV-Gali-2019}
The largest calculation contains a single NV$^{-}$ center in a 4096-site
diamond framework, corresponding to 4095 atoms after defect formation,
and approaches the dilute-defect limit.

The remainder of the manuscript is organized as follows.
Section~\ref{sec:methodology} develops the periodic velocity-gauge
formalism, the current operator including nonlocal pseudopotential
contributions, density-matrix propagation, the collinear spin extension,
and the formal current-based optical response.
Section~\ref{sec:implementation} describes the current numerical
realization in \texttt{RMG}, including implementation-specific choices
in the perturbation, current evaluation, propagation, parallel
organization, and spectral processing.  Section~\ref{sec:results}
presents the H$_2$, hBN, Si, and \ch{CrI3} validation benchmarks.
Section~\ref{sec:nv} presents the NV$^{-}$-center application, including
selection-rule tests, active-space dependence, concentration effects,
the 4095-atom single-defect calculation, and collinear magnetic
configuration comparisons.
